\textbf{Learned microarchitectural controllers.} A growing body of work replaces hand-tuned heuristics with online or offline-trained predictors across every major structure Bouncer targets. RL memory schedulers \citep{ipek2008selfoptimizing} and RL prefetchers \citep{bera2021pythia} learn performance-coupled policies online; learning-based cache replacement imitates Belady's optimum with PC-indexed predictors \citep{jain2016hawkeye}, distilled neural models \citep{shi2019glider}, or signature-based hit predictors \citep{wu2011ship}; and neural and geometric-history branch predictors \citep{jimenez2001perceptron, seznec2006tage} have become the dominant adaptive prediction substrate. These designs share a failure surface: their advantage over a simple baseline (LRU/RRIP, a stride prefetcher, a bimodal predictor) is workload- and phase-dependent and can silently invert under distribution shift. Bouncer is agnostic to which of these controllers it wraps \emph{among the set-local, reseed-identifiable ones it can faithfully audit}; rather than proposing a new policy, it provides, for that class, the missing runtime guarantee that the controller costs, in expectation, no more than a small bounded amount over its own known-safe fallback.

\textbf{Set dueling and adaptive policy selection.} Set dueling \citep{qureshi2007dip} dedicates a handful of sampled sets to each of two competing policies and lets a single saturating counter steer the followers, a mechanism extended to re-reference prediction \citep{jaleel2010rrip}, shared-cache thread-aware insertion \citep{jaleel2008tadip}, and evicted-address filtering \citep{seshadri2012eaf}, and rooted in the sampling-monitor methodology of utility-based partitioning \citep{qureshi2006ucp} and sampling dead-block prediction \citep{khan2010sdbp}; a recent retrospective documents its industrial adoption and limits \citep{qureshi2007retro}. Bouncer borrows the dedicated-sets-plus-counter substrate but repurposes it from policy \emph{selection} to policy \emph{trust}: where DRRIP duels two heuristics to pick the empirical winner, Bouncer duels the learned controller against its fallback to \emph{certify} a competence advantage and, crucially, keeps the run-C/run-fallback assignment \emph{secret}. Classic set dueling assumes cooperative policies and is content to follow whichever wins; Bouncer must instead defend against a controller that may be steered or drifting, which is why its dueling sets feed a change detector and a trust gate rather than a preference counter.

\textbf{Safe RL, runtime assurance, and Simplex.} The closest architectural ancestor is Simplex \citep{sha2001simplicity}: a simple, formally verified safety controller bounds an unverifiable high-performance one via a decision module. Shielding \citep{alshiekh2018shielding} synthesizes a correct-by-construction shield from a temporal-logic specification, and constrained- and Lyapunov-based safe RL \citep{altman1999cmdp, achiam2017cpo, chow2018lyapunov, berkenkamp2017stability} keep policies inside a certified-safe set, with safe policy improvement \citep{thomas2015hcpi, laroche2019spibb} reverting to the data-collecting baseline wherever improvement is not statistically supported (surveyed in \citep{garcia2015survey}). Bouncer inherits the ``use simplicity to control complexity'' posture and the bootstrap-to-baseline reflex, but departs in what it must prove. Simplex and shielding require a \emph{verified model of the safety envelope}; constrained RL requires a known cost function. Bouncer needs neither: microarchitectural ``safety'' is not a state-space invariant but a relative performance floor, so it certifies only that the learned controller's \emph{realized reward} beats the fallback's on randomized dueling sets, with a bounded-regret safety floor (our Lemma) that holds without any reachability proof or verified plant model. The closest relaxation of Simplex's model requirement, Black-Box Simplex \citep{mehmood2022blackbox}, still forward-simulates a reachability/admissibility check and defines safety as a state-space invariant; Bouncer instead defines it as a measured relative-competence floor and needs no plant model at all. Concurrent runtime-assurance work certifies recovery \emph{deadlines} for adapting controllers via conformal prediction \citep{shojaei2026conformal}; Bouncer is complementary, certifying a competence floor rather than a recovery time, and adds the secrecy-based mimicry resistance that an observable certificate lacks.

\textbf{Change detection and OOD monitoring.} Tier-B runs a one-sided CUSUM
\citep{page1954cusum}, using quickest-detection theory
\citep{lorden1971reacting,moustakides1986optimal} and sequential threshold sizing
\citep{siegmund1985sequential,tartakovsky2014sequential}. Tier-A instead draws on
input monitoring: confidence baselines \citep{hendrycks2017baseline}, conformal
validity \citep{vovk2005algorithmic,lei2018distribution}, concept-drift detection
\citep{gama2014survey}, and representation-level shift tests
\citep{rabanser2019failing}. Such monitors measure whether inputs move; Bouncer's
duel measures whether realized controller reward falls relative to a fallback.

Related label-free deterioration tests include the Suitability Filter
\citep{pouget2024suitability}, sequential harmful-shift detection
\citep{amoukou2024sequential}, and D3M \citep{d3m2024}. Idealized D3M provides the
exact false-alarm guarantee; its practical variant approximately inherits it.
Bouncer's distinction is a secret, reseeded, hardware-oriented competence sample
with an expected-regret gate. The illustrative storage count is 406 bytes, but
the proposed quantization and RTL timing are unvalidated. Its attribution also
requires representative sampling, set-locality, and reseed-identifiability.

\textbf{Microarchitectural security and co-tenant attacks.} The controllers Bouncer audits are also attack surfaces. Contention and timing channels on caches \citep{osvik2006cache, yarom2014flushreload, liu2015llc, gruss2016flushflush} and occupancy channels \citep{shusterman2019occupancy} leak across cores and VMs in shared clouds \citep{ristenpart2009heyyou}, while branch predictors \citep{evtyushkin2016jumpoveraslr, evtyushkin2018branchscope} and prefetchers \citep{gruss2016prefetch, guo2022adversarialprefetch} expose secret adaptive state that an attacker can read out or steer. A learned controller widens this surface: its policy can be poisoned into amplifying leakage or into co-tenant denial of service. Prior defenses harden specific channels; Bouncer is orthogonal and complementary, bounding the \emph{aggregate} damage any steered controller can do by capping its cost at the vetted fallback. Its mimicry resistance (our Proposition) follows precisely from the secrecy of the set assignment, so an attacker who can observe and game stealthy counter-based detectors \citep{gruss2016flushflush} still cannot tell which accesses are being scored.

\textbf{ML-for-systems robustness.} Deployed learned components are brittle in exactly the regimes Bouncer guards. Oracle-imitation cache controllers lose accuracy off-distribution \citep{liu2020imitationcache}, learned indexes degrade on irregular data \citep{marcus2020benchmarklearnedindex}, and the Park benchmark documents why RL-for-systems policies behave unlike RL-for-games under noisy, non-stationary workloads \citep{mao2019park}; this is the hidden, system-level risk catalogued in \citep{sculley2015technicaldebt}. Production systems respond with bespoke guardrails: SLO-clamped colocation \citep{lo2015heracles}, OOM-bounded autoscaling \citep{rzadca2020autopilot}, and fault-triggered fallback in learned software \citep{carstens2023smartchoices}. Bouncer generalizes these one-off, fault- or threshold-triggered safeguards into a single hardware-cheap mechanism with a \emph{continuous} competence audit and a provable performance floor, applicable to any \emph{set-local, reseed-identifiable} learned microarchitectural controller rather than one tuned guardrail per deployment.
